\section{Measurement and evaluation methods}
\label{sec:methods}

\subsection{Estimands and censoring}

We define five primary estimands before implementation. (1)~Terminal
scientific failure is the event that the trajectory's final scientific
claim does not satisfy the independent adjudicator's validity criteria,
regardless of whether execution completed. (2)~Failure onset is the first
trajectory step where the adjudicator identifies a state that, if left
uncorrected, would cause terminal scientific invalidity; onset is measured
as interval-censored between the last verified-valid step and the first
detected-invalid step. (3)~Recovery is an intervention after failure onset
that results in subsequent independently verified scientific validity, with
the fault removed, downstream state consistent, and no new validity
violations introduced. (4)~Warning lead time is the normalized trajectory
horizon between a warning signal and the later of terminal check detection
or adjudicated onset. (5)~Attribution is a hierarchical label drawn from
\{known single mechanism, known composition, unknown, ambiguous,
insufficient evidence\}; the unknown and ambiguous categories are retained
in evaluation rather than excluded.

Trajectories that do not complete are classified into eight censoring
categories: policy termination, budget exhaustion, infrastructure failure,
injector failure, observer failure, verifier failure, scientific task
failure, and unknown. Intention-to-treat analysis retains all valid
assignments; a per-protocol sensitivity excludes injector and
infrastructure failures but does not replace the primary analysis.

\subsection{Hierarchical inference}

The task instance is the principal independent sampling unit. Repeated
agent seeds within a task estimate conditional variability but do not
inflate the independent sample count. Confidence intervals use
task-grouped bootstrap resampling: task templates are resampled first, task
instances within template second, and all trajectories or repeated seeds
within instance are retained or resampled as specified. At least 2{,}000
bootstrap replicates are used for final intervals. Model-based intervals
from hierarchical mixed models are reported as a sensitivity comparison.

Multiplicity across secondary contrasts (per workload, per fault family, per
baseline) is controlled with Benjamini-Hochberg false-discovery-rate
adjustment within predeclared families. Exploratory analyses are labeled as
such and do not inherit confirmatory language.

\subsection{Failure-dynamics analysis}
\label{sec:methods_failure_dynamics}

The failure-dynamics analysis (experiment R1-E1) estimates hazard and
propagation from the confirmatory benchmark. Exposure variables include
trajectory age (normalized by remaining horizon), fault family, severity,
persistence, injection band, accumulated unresolved state divergence,
evidence freshness, and model/workload configuration. Outcomes include
adjudicated failure-onset interval, terminal execution failure, terminal
scientific invalidity, and propagation edges between typed subsystem states.

We use Kaplan-Meier curves for description and a discrete-time mixed-effects
hazard model with complementary log-log link for inference. Task instance
and template random effects capture repeated measures; model, workload, and
fault family enter as fixed effects with prespecified interactions.
Interval-censoring sensitivity relaxes the exact-onset assumption, and
clustered bootstrap intervals over task templates provide distribution-free
uncertainty. Competing-risk sensitivity distinguishes policy termination
from scientific failure.

A preregistered failure law—for example, that the hazard increases
monotonically with accumulated unresolved state divergence after adjusting
for fault severity—is tested for replication across at least two workload
families after removal of terminal-adjacent events and under an alternative
onset rule. The law is confirmed only if its direction, practical magnitude,
and grouped interval survive all sensitivity conditions.

\subsection{Warning and attribution baselines}
\label{sec:methods_warning}

The warning analysis (experiments R1-E2 and R1-E3) compares seven model
configurations: (1)~a terminal-verifier-only baseline that declares failure
only after the adjudicator's final check; (2)~runtime-rule alarms triggered
by predefined error patterns; (3)~agent self-confidence extracted from the
agent's own output; (4)~a text-only trace critic operating on raw log text;
(5)~a typed-event temporal model that consumes structured event histories
without state differences; (6)~the full execution-grounded model that
additionally consumes state differences and evidence provenance; and
(7)~oracle fault metadata, reported only as an unattainable upper bound
using hidden injection truth.

The primary endpoint for H1 is median normalized warning lead time among
terminally invalid trajectories at an operating point with no more than
5\% false warnings on no-fault valid trajectories. The operating threshold
is selected on a dedicated calibration split, never on the test set. A
warning counts only if it precedes both the first terminal-only signal and
the minimum declared reaction window.

Secondary endpoints include time-dependent AUPRC and AUROC, integrated Brier
score and log loss, calibration intercept and slope with reliability plots,
risk-coverage and abstention-utility curves, lead-time distribution by fault
and workload, and alert burden per successful trajectory. Post-shift
performance (after environment or provider change) is measured as a
robustness check.

Attribution uses hierarchical classification over the five label categories.
Metrics include macro-F1 over known top-level families, top-$k$ recall,
hierarchical distance, multilabel set F1 for compositions, and
unknown-family detection rate. Results are reported both including and
excluding ambiguous cases. The cause distribution is calibrated on the same
split as the hazard model, using isotonic or temperature scaling as
appropriate.

The model architecture is deliberately compact: a discrete-time hazard
regression or small temporal encoder with post-hoc calibration. Model
complexity is expanded only if the simple baseline fails for a diagnosed
capacity reason. The primary scientific contribution is the measurement
artifact and the empirical result, not a novel neural architecture.
